\begin{lemma}[A bad multi-sequence has a nontrivial deciding front]
Let $r$ be a reflexive relation on $Q$.  If
$h:\mathsf{IncSeq}\to Q$ is locally constant and bad for $r$, then there is
a super-sequence
\[
f:F^h\to Q,
\qquad F^h=\mathsf{DecidingFrontSet}(h),
\]
which is bad, and $\emptyset\notin F^h$.
\end{lemma}
\begin{proof}
By \noderef{DecidingFront}, the set
$F^h=\mathsf{DecidingFrontSet}(h)$ is a front on $\mathbb N$.  By
\noderef{DecidingValue}, there is a map
$v:F^h\to Q$ satisfying
\[
v(s)=h(x)
\quad\text{whenever }s\text{ is a proper prefix of the range of }x.
\]
Package this front and value map as
\[
f=\langle \noderef{DecidingFront}(h),v\rangle,
\]
which has the required \noderef{SuperSequence} type.  The displayed value
equation is exactly the hypothesis required by \noderef{BadMultiToSuper};
therefore \noderef{BadMultiToSuper}, together with the assumed
\noderef{BadMultiSequence} badness of $h$, proves
$\mathsf{BadSuperSequence}(r,f)$.

It remains to prove that the front is nontrivial.  Suppose instead that
$\emptyset\in F^h$.  Unfolding \noderef{DecidingFrontSet}, this gives
$\mathsf{DecidingPrefix}(h,\emptyset)$.  For every increasing enumeration
$x$, the empty set is a proper prefix of its range: use the first value
$x(0)$ as the witnessing next element, and use monotonicity of $x$ to show
that no element of its range is below $x(0)$.  In particular this holds for
$\mathsf{IncSeqId}$ and for $\mathsf{ShiftMap}(\mathsf{IncSeqId})$, where
the relevant increasing enumerations are supplied by
\noderef{IncSeqId} and \noderef{ShiftMap}.  The deciding-prefix equation
therefore gives
\[
h(\mathsf{IncSeqId})
=h(\mathsf{ShiftMap}(\mathsf{IncSeqId})).
\]
Reflexivity of $r$ now gives the relation between these two values, while
\noderef{BadMultiSequence} gives its negation at
$\mathsf{IncSeqId}$.  This contradiction proves
$\emptyset\notin\mathsf{DecidingFrontSet}(h)$, completing the lemma.
\end{proof}
